The unit-testing performed, throughout the development process, 
focused on public methods of the ,``LidDrivenCavity" class (LDC), and treated the ,``SolverCG'' class (SCG),
as a unit onto itself. The unit tests kept the same checking conditions when extending the original
serial code with MPI and OpenMP, however significant refactoring was necessary to facilitate these equivalent checks.
The source code for these tests can be found in \app\ref{Appendix::Test}.
\subsection{``LidDrivenCavity.h"}
The tests written for the LDC were deigned to validate the public methods the user would be accessing within the main script.
The setters and getters such as ``SetDomainSize'' were not tested due to their simplistic nature.
If more complexity was used in the afformentioned methods then they also would have been tested. 
\subsubsection{LidDrivenCavity::PrintConfiguration()}
The ``PrintConfiguration'' test utilised a stream-buffer as shown in \code\ref{lst:LidDrivenCavityTest}. After rerouting the output stream, the printed information
is extracted using a pattern search, via a regular expression. Finally, a vector is stored contained the printed values and the expected value.
Using the ``compareSolution()'' defined in \code\ref{lst:MyTest} the 2-norm or the error is computed and checked against a tolerance of 1e-3.
\subsubsection{LidDrivenCavity::WriteSolution()}
On the other hand the ``WriteSolution'' test involved conditions that were  agnostic of the accuracy of the SCG class.
The conditions require that the solver has successfully written out to the file requested through a comprehensive list of checks. 
These range from the file existing; the maximum number of values on a given line being the six paramters: 
x co-ordinate, y co-ordinate, vorticity, stream function, x velocity and y velocity;
and finally that there are the correct number of lines / grid points being printed.
\subsection{``SolverCG.h"}
Whilst the previous unit tests weren't concerned with the accuracy of the data being calculated, the following tests validate the ``SolverCG::Solve'' method.
Using the below equation \eq\ref{eq:nabla} that spatial Poisson problem was defined. 
\begin{equation}
    - {\nabla}^2 \psi \hat{\mathbf{k}} = \omega
    \label{eq:nabla}
\end{equation}
The main way the test changed once MPI had been implemented, was that the local error between the numerical and analytical solutions had to be precomputed and {\it reduced} to every other rank.
The tolerance used in the following tests was equivalent to the one set in the main solver file to maintain consistency.
\subsubsection{SolverCG::Solve() - Analytical}
Using a sinusoidal solution allowed a varying frequency in the stream function, shown in \eq\ref{eq:psi}.
The significance of this function was that the boundary conditions at the walls of the cavity were enforced given that at $x=L_x$, $y=L_y$ and at the origin.
\begin{equation}
    \psi = \sin\bigl(\pi * k * x\bigr)\sin\bigl(\pi * l * y\bigr),\quad k = \frac{2\pi}{L_x},\quad l = \frac{2\pi}{L_y}
    \label{eq:psi}
\end{equation}
Combining \eqs \ref{eq:nabla} and \ref{eq:psi}, gave the input solution to the ``Solve'' method as it used a conjugate gradient algorithm to solve for $x$ in $Ax=B$.
This input is found below in \eq\ref{eq:vort}.
\begin{equation}
    \omega = \pi^2(k^2+l^2)\sin\bigl(\pi * k * x\bigr)\sin\bigl(\pi * l * y\bigr),\quad k = \frac{2\pi}{L_x},\quad l = \frac{2\pi}{L_y}
    \label{eq:vort}
\end{equation}
As seen in \code\ref{lst:SolverCGTest} the domain used was of size 6,000,000 grid points and could be run with a variable number of MPI ranks as shown in \code\ref{lst:MyTest} through the Boost global fixture.
\subsubsection{SolverCG::Solve() - Zero Input}
Constrastingly, the ``Zero Input'' test was used to ensure that during the development of the parallelisation, the solver still gave back a zero vector.
This functionality  was implemented to avoid errors from failing to converge with such a small vorticity input.
